\section{Introduction: The In-Situ Resource Utilization Imperative in Lunar Life Support}
\label{sec:introduction}

Establishing an enduring human presence on the lunar surface requires transitioning from open-loop, resupply-dependent life support architectures to closed-loop, regenerative Biogenerative Life Support Systems (BLSS)\cite{wheeler2017agriculture,paul2022plants,ming1989lunar}. Terrestrial logistics impose an unsustainable launch mass burden for multi-year surface missions, where crew metabolic requirements for water, oxygen, and fresh nutrition scale linearly with crew complement and mission duration\cite{hanford2004advanced,wheeler2023space}. Higher-plant bioregenerative architectures address this vulnerability by using photoautotrophic crop canopies to scrub atmospheric carbon dioxide ($\mathrm{CO}_2$), generate metabolic oxygen ($\mathrm{O}_2$), purify greywater via foliar transpiration, and produce fresh, nutrient-dense biomass\cite{wheeler2017agriculture,ming1989lunar,wheeler2023space}.

Historically, space mission trade studies evaluated two distinct edaphic paradigms for planetary crop production: importing prefabricated hydroponic hardware from Earth, or cultivating crops directly in untreated, in-situ planetary regolith\cite{paul2022plants,ming1989lunar,wamelink2019crop}. Importing inert root-support matrices (such as mineral rockwool, polyurethane foams, or expanded clay) and prefabricated polymeric fluid channels from Earth imposes launch-mass penalties that severely restrict the deployable agricultural footprint during early outpost establishment\cite{ming1989lunar,hanford2004advanced,wamelink2019crop}. Conversely, direct edaphic cultivation in raw lunar regolith---often termed ``direct regolith agriculture''---operates under the premise that lunar dust can serve as a functional analog to terrestrial soil once augmented with water and soluble fertilizers\cite{ming1989lunar,wamelink2019crop,atkin2024bioremediation}.

However, pioneering botanical experiments utilizing returned Apollo specimens and high-fidelity lunar simulants demonstrate that raw regolith imposes severe physiological penalties\cite{paul2022plants,ming1989lunar}. Planetary-scale exposure to unattenuated solar wind protons, galactic cosmic rays, and hypervelocity micrometeorite comminution has produced a geochemically hostile matrix characterized by submicroscopic zero-valent metallic iron ($\mathrm{npFe}^0$), unweathered razor-sharp mineral edges, and volatile abiotic surface reactivity\cite{paul2022plants,ming1989lunar,hendrix2024reactivity,turci2015mechanistic}. Furthermore, single-particle wettability determinations on returned Chang'e-5 samples reveal pronounced intrinsic hydrophobicity across pristine lunar dust grains\cite{zhang2025wettability}. In parallel, the reduction of gravitational acceleration to one-sixth terrestrial gravity ($1/6\,g$) drastically lowers the dimensionless Bond number, causing capillary retention forces to dominate over gravitational drainage\cite{ming1989lunar,steinberg2007fluid,jones1999microgravity}. As demonstrated by recent suborbital centrifuge flight experiments simulating lunar gravity ($1.76\,\mathrm{m/s}^2$), fluid menisci exhibit rounded curvatures and boundary film thickening, while hydrophobic surface dust severely impedes water infiltration and liquid-gas equilibration\cite{hasenstein2026assessing}. Mitigating root-zone hypoxia under fractional gravity requires coarse, well-aerated aggregates ($0.5\text{--}2.0\,\mathrm{mm}$)\cite{stout2001evidence,hasenstein2026assessing}; however, core stratigraphy from Apollo and Chang'e-5 drill cores confirms that such millimeter-scale fractions dominate only at depths of several meters\cite{mckay1991lunar,zhao2023variations}, imposing prohibitive mechanical excavation and sifting burdens on early human outposts.

To overcome both the launch-mass penalty of imported hardware and the biological toxicity of raw regolith, life-support engineering is converging toward an in-situ bio-infrastructure paradigm\cite{ming1989lunar,taylor2005microwave,fateri2019solar}. Rather than attempting to adapt plant genetics to an inherently hostile mineral substrate, In-Situ Resource Utilization (ISRU) thermal consolidation processes can transform raw regolith into engineered porous ceramics\cite{ming1989lunar,taylor2005microwave,xi2026multiphysics}. High-temperature sintering passivates reactive geochemical interfaces, rounds abrasive grain morphology, and enables the precision fabrication of hydroponic troughs, multiphase fluid-delivery wicks, and mycoponic scaffolds\cite{ming1989lunar,xi2026multiphysics,goulas2017process}. This review analyzes the physicochemical, toxicological, fluid-mechanical, and systems-engineering principles that govern the transition from raw regolith farming to sintered ceramic bio-infrastructure for sustained lunar settlement.
